\subsection{Frontend and normalised clauses}
\label{sec:frontend}

The frontend accepts OWL functional syntax directly.  OWL/XML and RDF/XML are
parsed with the Horned-OWL Rust library \cite{lord2024horned} and Turtle with a
separate RDF parser; both paths produce a functional-syntax handoff into the
same normaliser.  RDF input must map completely onto OWL structures; leftover
triples or disconnected blank nodes produce an unsupported result instead of a
partial ontology.  The classifier performs no network access, so an import
closure must be merged into one document beforehand.  The frontend
also checks the OWL~2~DL global restrictions and refuses an ontology that
violates them with an explanation and a dedicated exit code, so that a
malformed input is reported as such instead of being classified under a
repaired reading.

A streaming first pass resolves IRIs, builds a typed source representation,
and records the source profile that drives routing: counts of logical TBox,
RBox, ABox, rule, declaration, and annotation axioms; declared and distinct
classes, properties, and individuals; counts of each axiom and concept
constructor, including conjunction, disjunction, complement, existential and
universal restrictions, cardinalities, nominals, value and self restrictions,
and datatype use; occurrences of the bottom class and bottom properties; and
the maximum concept depth, role-chain length, and cardinality bound.  The pass
also computes a Konclude-compatible expressivity code by porting Konclude's
expressiveness calculation, including its preprocessing effects for active
transitive roles, inverse partners, and reachable nominal domains; on the 592
benchmark ontologies every code matches the official Konclude binary.  The same
pass evaluates several source-level certificates that the router uses, for
example that a positive ABox cannot change any named-class subsumption, or
that no number-restricted role shares a hierarchy component with an inverse,
chained, or transitive role.

Normalisation produces DL clauses $\Gamma\rightarrow\Delta$ whose body
$\Gamma$ is a conjunction and whose head $\Delta$ is a disjunction of atoms.
Atoms are concept atoms $A(t)$, role atoms $R(s,t)$, and equality atoms
$s\approx t$; an inequality is an equality atom in a body.  Terms are the
central variable $x$, its neighbour variables, named individuals, auxiliary
constants, and function terms $f(t)$ standing for existential witnesses.  Atom
lists are sorted and de-duplicated so that structurally equal clauses have one
representation.  The clausifier follows the structural transformation used by
hypertableau and consequence-based reasoners
\cite{motik2009hypertableau,bate2018jair}: axioms are converted to negation
normal form; each structurally distinct complex subconcept receives one fresh
definer concept shared across its occurrences; a positive existential
restriction $\exists R.C$ becomes the clauses $\Gamma\rightarrow R(x,f(x))$ and
$\Gamma\rightarrow C(f(x))$ with a fresh function symbol $f$; a universal
restriction receives introduction clauses only where it occurs at negative
polarity; qualified number restrictions become at-least and at-most clauses
whose heads contain equalities between witnesses, and are additionally
recorded as first-class cardinality definitions for the cardinality routes; a
nominal $\{o\}$ becomes a fresh proxy concept with a recorded mapping; inverse
roles produce bridge clauses $R(x,y)\rightarrow S(y,x)$ and
$S(x,y)\rightarrow R(y,x)$; and ABox assertions become ground clauses.  An
augmentation pass adds transitivity, role-chain, domain, and range clauses,
using role automata for chains, including chains with inverse members.

Data ranges and literals are abstracted to concepts over data nodes.  A
datatype oracle decides membership, equality, and inclusion between the data
concepts that occur in the clause set according to the OWL~2 datatype map and
emits an ordinary concept clause for each decided relation.  When a decision
procedure returns unknown it emits nothing, so an unsupported datatype corner
degrades to the sound abstraction.  Datatype reasoning therefore stays in the
frontend, and no calculus rule or proof is specific to datatypes.  Side
structures carry what a bare clause vector cannot: the internal-name-to-IRI
map, named and declared classes, whether the RBox is EL-safe, ABox facts
established during parsing, nominal and ABox provenance with a completeness
flag, cardinality definitions, definer provenance, the normalised source TBox,
DL-safe rules, the source profile, and the selected route.  A worker that
receives clauses but would lose typed ABox or cardinality metadata must
decline.  Large sources stream clauses to a file and side data to a meta file,
so the clause set is never copied through an intermediate runtime or held
twice in memory.

\subsection{\EL{} completion}
\label{sec:elc}

The EL worker classifies ontologies in the \EL{} fragment without nominals and
concrete domains.  It implements the completion procedure of Baader, Brandt,
and Lutz \cite{baader2005el} with the indexing and scheduling of ELK
\cite{kazakov2011concurrent,kazakov2014elk}, over the frontend's clauses.  Its
output is the full set of subsumers of every class (KM computes no transitive
reduction), the unsatisfiable classes, and the consistency of the ontology.
Figure~\ref{fig:flow-el} shows the procedure and
Table~\ref{tab:opt-el} in the appendix lists every optimisation with its
default and origin.

The worker first maps each clause onto one of seven normal forms,
$A\sqsubseteq B$, $A\sqcap B\sqsubseteq C$, $A\sqsubseteq\exists R.B$,
$\exists R.A\sqsubseteq B$, $A\sqsubseteq\bot$, $R\sqsubseteq S$, and
$R\circ S\sqsubseteq T$, plus reflexive roles, with all names interned to dense
integers.  An $n$-ary conjunction becomes a left-deep chain of binary
conjunctions whose auxiliary concepts are shared between identical prefixes; a
domain axiom becomes $\exists R.\top\sqsubseteq B$; and the two Skolem
half-clauses $\Gamma\rightarrow R(x,f(x))$ and $\Gamma\rightarrow C(f(x))$ are
reassembled into $A\sqsubseteq\exists R.C$.  Every clause that matches none of
these shapes, such as a disjunctive head, an equality, a range, or an inverse
bridge, goes to a residual list.  Under the automatic route a zero-copy screen
that mirrors this mapping runs in the frontend, so an ontology whose residual
would be non-empty never reaches the EL route.

If the residual is empty and the ontology consists only of acyclic named
subclass edges, the worker computes the closure in reverse topological order
\cite{kahn1962topological} and returns without saturation; a cycle sends it to
the general procedure.  Otherwise it builds a rule index that stores, for every
trigger concept, one contiguous record of its subclass, conjunction,
existential, and existential-left-side axioms, precomputes the
reflexive-transitive closure of the role hierarchy, seeds every class with
itself and $\top$, and saturates.  Saturation keeps, for every class, a label
of derived subsumers, forward edges, backward links keyed by target and exact
role, and a propagation store for left-side existentials.  A new subsumer $d$
of $c$ fires, in this order, the subclass rule, the conjunction rule joined
from the smaller side, the bottom rule, the existential rule that adds an edge
$(c,R,F)$, the backward propagation of $\bot$ to all predecessors, and the
registration of every axiom $\exists R.d\sqsubseteq e$ in the propagation store
together with its application to existing predecessors.  A new edge
$(c,r,d)$ fires the stored propagations of $(d,r)$, propagates $\bot$ from $d$,
applies role chains in both directions, and lifts the edge to every super-role,
so that the existential and chain rules match roles exactly.  Each two-premise
rule fires from whichever premise arrives second, so the fixpoint is the same
under any schedule.  Items are processed per context: a context is drained
completely before the next is activated, which keeps its tables in cache.
A class is unsatisfiable when $\bot$ is in its label, and the ontology is
inconsistent when $\bot$ is in the label of $\top$.

With the certificate off, which is the default and the setting of the
\texttt{elc} route, a non-empty residual makes the worker decline before any
saturation.  Two certificate modes extend the procedure beyond pure \EL{}.
Both read the saturated state as an interpretation $\mathcal{I}$ whose elements
are the satisfiable class nodes, with $x_C\in D^{\mathcal{I}}$ exactly when $D$
is in the label of $C$ and role extensions given by the derived edges; at the
fixpoint $\mathcal{I}$ satisfies every \EL{} clause.  The check mode compiles
each residual clause into a constraint over $\mathcal{I}$, with every Skolem
function pinned to its own witness node so that distinct functions denote
distinct elements, and evaluates it by a bounded join.  If every residual
clause holds, $\mathcal{I}$ is a model of the whole ontology, every entailed
subsumption holds in it, and the \EL{} answer is exact.  The repair mode, used
by the certified EL route and by the production portfolio, treats the \EL{}
answer as a lower bound and builds upper bounds by repair, in the spirit of
lower- and upper-bound reasoning \cite{zhou2015pagoda}: each of two to eight
passes forks the saturated state, makes one head atom of every violated
residual instance true, re-saturates, and repeats until every residual clause
holds.  A choice that makes a satisfiable named class unsatisfiable is banned
and the pass restarts, as in conflict-driven search
\cite{marquessilva1999grasp}.  A class whose additional named subsumers in the
repaired models have an empty intersection is answered exactly; if at most 64
classes remain undecided, the worker publishes the rest and the
consequence-based engine classifies the undecided classes, and otherwise it
declines.

Two further paths change only how the computation is scheduled.  Context-parallel
saturation partitions the contexts statically over workers that exchange edge
halves in batched messages and detect termination by credit counting
\cite{mattern1989credit}; it rebuilds labels in sorted order so that its output
is byte-identical to the serial run, and the router enables it only for large
ABox-free \EL{} terminologies.  A source-level scanner classifies very large
flat taxonomies before the frontend runs, and small structured sources are
classified inside the supervisor process to avoid a serialisation round trip.

\begin{figure}[t]
\centering
\begin{tikzpicture}[flow, node distance=3.2mm and 9mm]
\node[step=60mm] (a) {normalised DL clauses};
\node[step=60mm, below=of a] (b) {map to NF1--NF7 and reflexive roles; collect residual};
\node[test=60mm, below=of b] (c) {residual empty and acyclic subclass edges only?};
\node[step=34mm, right=of c] (c1) {topological closure};
\node[test=60mm, below=of c] (d) {residual present and certificate off?};
\node[term=34mm, right=of d] (x) {decline};
\node[step=60mm, below=of d] (e) {rule index, role closure; seed $C\sqsubseteq C$, $C\sqsubseteq\top$};
\node[step=60mm, below=of e] (f) {saturate: context worklist, backward links, propagations; serial or context-parallel};
\node[test=60mm, below=of f] (g) {residual present?};
\node[step=34mm, right=of g] (h) {check or repair the canonical model; residue to CB};
\node[term=80mm, below=7mm of g, xshift=10mm] (out) {subsumers, unsatisfiable classes, consistency};
\draw[->] (a) -- (b); \draw[->] (b) -- (c);
\draw[->] (c) -- node[lab,above]{yes} (c1);
\draw[->] (c1.east) -- ++(0.3,0) |- (out.east);
\draw[->] (c) -- node[lab,right]{no} (d);
\draw[->] (d) -- node[lab,above]{yes} (x);
\draw[->] (d) -- node[lab,right]{no} (e);
\draw[->] (e) -- (f); \draw[->] (f) -- (g);
\draw[->] (g) -- node[lab,above]{yes} (h);
\draw[->] (h) -- node[lab,right]{fail} (x);
\draw[->] (h.south) |- node[lab,pos=0.3,right]{pass} ($(out.north east)+(0,0)$);
\draw[->] (g) -- node[lab,right]{no} (out);
\end{tikzpicture}
\caption{\EL{} completion.  The worker declines whenever a clause falls
outside the normal forms and no certificate is requested; a certificate either
proves that the canonical model satisfies the residual clauses, answers all
but a bounded residue that the consequence-based engine classifies, or
declines.}
\label{fig:flow-el}
\end{figure}

\subsection{Disjunctive consequence-based engine}
\label{sec:cb}

The consequence-based (CB) engine implements the disjunctive context calculus
of Bate et al.\ \cite{bate2016sriq,bate2018jair} with the nominal extension of
Tena Cucala, Cuenca Grau, and Horrocks \cite{tenacucala2018ijcai}, and it
follows the term order, literal ordering, and index design of the Sequoia
reasoner \cite{tenacucala2019sequoia,tenacucala2021payg}.  It is the
procedure of last resort for every ontology outside the EL fragment, and the
nominal-aware variant is the exact fallback of every specialist route.
Figure~\ref{fig:flow-cb} shows the procedure, and
Tables~\ref{tab:rules-cb} and~\ref{tab:opt-cb} in the appendix give the rules and
optimisations.

The worker interns the frontend's clauses into an ordered term representation
in which the variables, the neighbour variables, individuals, and function
terms $f(x)$ and $f(o)$ are 32-bit integers whose numeric order is the term
order, so function terms are maximal.  Body equalities move to the head as
inequalities, and a clause whose role body atom avoids the central variable
must already have been rewritten by the frontend, which compiles transitivity
and role chains into fresh propagation concepts.  Individuals survive only in
nominal mode; outside it, a clause that mentions one is dropped and counted.
Two frontend options shape the clause set on the production route:
polarity-gated definitional clausification, which emits only the direction of
each definer that its occurrences need \cite{plaisted1986structure} and thereby
removes unguarded excluded-middle disjunctions, and triggered absorption, which
compiles triggerable antecedents into guarded clause bodies
\cite{steigmiller2014absorption}.  The worker then prepares, once per
classification, an immutable ontology shared by all threads: clause indexes by
body predicate, a guarded-role index for the absorbed shape
$C(y)\wedge R(x,y)\rightarrow D(x)$, trigger bits for the Succ rule, a quotient of
asserted individual equalities, and recognised nominal enumerations
$A\equiv\{o_1,\ldots,o_n\}$.

Classification creates one root context with core $\{A(x)\}$ for every named
query class $A$ and one for $\top$.  A context holds clauses
$\Gamma\rightarrow\Delta$ over its central variable and neighbour terms and
runs a given-clause loop: it pops a clause, discards it if an active clause
subsumes it, and applies the inference rules to the head literals that the
ordering makes maximal.  Hyper resolves an ontology clause with one maximal
head literal of each premise under a substitution that keeps the central
variable fixed, except in the ground context, where it may bind individuals.
Eq paramodulates with maximal equalities, oriented with the larger term first;
Ineq removes reflexive inequalities; and Factor merges two equalities with a
common left side, which together handle functional roles and number
restrictions in the manner of ordered paramodulation
\cite{bachmair1994rewrite}.  Elim discards tautologies, duplicates, and clauses
subsumed by an active clause, and removes active clauses that a new clause
subsumes.  When a context derives a maximal trigger literal on $f(x)$, Succ
creates or reuses the successor context for $f$ and installs the trigger there
as a hypothesis $p\rightarrow p$; Pred returns to the predecessor every
function-free clause whose body consists of hypotheses it was sent, with the
successor's core added to the body.  In nominal mode a single ground context
receives every assertion and every top-level fact instantiated at every
individual, and the rules r-Succ, r-Pred, Join, and Nom of the nominal calculus
exchange ground atoms with it; Nom introduces fresh nominals for an at-most
restriction over a nominal neighbour in stable blocks, and a run that exceeds
its budget of 4,096 fresh nominals declares itself incomplete.

The successor schedule decides how many contexts are created.  The default
central strategy keys a successor by the image of its unit triggers and of its
conditionally derived single triggers, which follows the shape of the Succ
rule in the published calculus \cite{bate2016sriq} and the eager strategies
of the context framework \cite{simancik2014fpt}; every other trigger arrives as
a hypothesis.  The per-function strategy, which creates one empty-core context
per function symbol, is the schedule of the single-threaded retry.  The literal
ordering is a port of Sequoia's: predecessor-trigger atoms are smallest and
pairwise incomparable; distinct concepts on the same term are mutually
incomparable by default, so Hyper resolves on every disjunct; and when some
clause head contains a disjunction over an internal definer, the definers are
totally ordered above the named concepts, which restores ordered resolution on
the normaliser's definers.

Contexts communicate only through Succ and Pred messages.  The driver takes all
pending messages as one batch, applies each to its target, and then saturates
and propagates every touched context once; propagation is semi-naive, scanning
only clauses and edges that are new since the last round.  The loop stops when
no message is pending.  The terms are restricted to the central and neighbour
variables, individuals, and one level of function terms, contexts are keyed by
finite atom sets, and clauses are sets, so the fixpoint is reached; a global
cap of 25 million messages remains as a resource backstop and, if it fires,
marks the run incomplete.  The subsumers of $A$ are the named $B$ with a unit
clause $\top\rightarrow B(x)$ in the root context of $A$, $A$ is unsatisfiable
when that context derives the empty clause, and the ontology is inconsistent
when the $\top$ context or the ground context does.  Classes in one strongly
connected component of told subclass edges share a root context.  Threads
take batches of queries from a shared cursor and each runs an independent
engine over the shared prepared ontology; because the subsumers of one query
depend on no other query, the partition changes no result.  An incomplete run
exits with a distinct status and publishes nothing.

\begin{figure}[t]
\centering
\begin{tikzpicture}[flow, node distance=3.2mm and 9mm]
\node[step=64mm] (a) {intern clauses; drop and count unsupported terms};
\node[step=64mm, below=of a] (b) {prepare the shared ontology: indexes, triggers, equality quotient, ordering mode};
\node[step=64mm, below=of b] (c) {root context with core $\{A(x)\}$ per query; ground context in nominal mode};
\node[step=64mm, below=of c] (d) {given-clause loop on maximal literals: Hyper, Eq, Ineq, Factor, Elim, Join, Nom};
\node[step=64mm, below=of d] (e) {propagate: Succ to successors, Pred to predecessors, r-Succ, r-Pred};
\node[test=64mm, below=of e] (f) {messages pending?};
\node[step=32mm, right=of e] (g) {apply the batch; saturate each touched context};
\node[term=32mm, right=of f] (x) {cap or budget hit: incomplete};
\node[term=64mm, below=of f] (out) {read off units $\top\rightarrow B(x)$ and empty clauses};
\draw[->] (a) -- (b); \draw[->] (b) -- (c); \draw[->] (c) -- (d); \draw[->] (d) -- (e); \draw[->] (e) -- (f);
\draw[->] (f.east) -- ++(0.35,0) |- node[lab,pos=0.25,left]{yes} (g.west);
\draw[->] (g.north) |- (d.east);
\draw[->] (f.south east) -- node[lab,below]{cap} (x.south west);
\draw[->] (f) -- node[lab,right]{no} (out);
\end{tikzpicture}
\caption{The consequence-based engine.  Each context saturates locally; Succ
and Pred messages are applied in batched rounds until none is pending.  A run
that hits a resource backstop publishes nothing.}
\label{fig:flow-cb}
\end{figure}

\subsection{Hypertableau procedures}
\label{sec:ht}

KM contains two hypertableau procedures: a native one written for KM after
HermiT's design \cite{motik2009hypertableau,glimm2014hermit}, and a
function-by-function Rust port of Konclude's completion and classification
algorithms \cite{steigmiller2014konclude,steigmiller2016thesis}, entered through
a bridge.  Both read the output of a converter that turns DL clauses into
hypertableau clauses.  The converter's only logical transformation is reverse
Skolemisation: the pair $\Gamma\rightarrow R(x,f(x))$ and
$\Gamma\rightarrow C(f(x))$ becomes the existential $\Gamma\rightarrow\exists R.C$,
and a disjunctive Skolem head is left unfolded because folding it would
strengthen a disjunction into a conjunction.  It also parses the role
hierarchy, encodes transitivity by propagated markers for universal
restrictions \cite{horrocks1999transinv}, emits equality clauses for at-most
and functionality constraints, emits first-class at-least and at-most
definitions when a role-separation certificate permits them, eliminates pairs
of complementary definers whose folding is independently derivable, and fences
the combinations it cannot carry, such as nominals with inverse roles or
inverse roles with number restrictions.  It counts every dropped or fenced
clause, and every hypertableau route declines if any clause relevant to it is
dropped.  Figure~\ref{fig:flow-ht} shows the two
procedures, and Tables~\ref{tab:opt-ht} and~\ref{tab:opt-bridge} in the appendix
list their optimisations.

The native procedure decides satisfiability by a depth-first search over a
completion graph whose facts carry dependency sets, the sets of branch levels
that support them.  Each step drains a queue of new facts through
hyperresolution: a new concept fact re-matches only the clauses whose body
contains it, a new edge only the clauses with that role in the body, and a new
node the clauses with an empty body \cite{motik2007opt}.  A clause with an
empty head raises a clash, a single atom is asserted, an existential becomes
an obligation, and a disjunction becomes a pending choice.  The step then
merges every node that carries a nominal into one survivor, expands
existential and at-least obligations on nodes that are not blocked, and scans
the pending disjunctions, which are watched so that only the ones touched by
a change are re-examined.  A disjunction with one live disjunct is asserted;
otherwise the search branches.  It tries the disjuncts in order of past
failures, asserts the negation of each failed disjunct before trying the
next, and returns a clash whose dependency set excludes the current level
without trying the remaining disjuncts, which implements semantic branching and
dependency-directed backjumping \cite{horrocks1999optimizing}.  When no choice
remains, first-class at-most restrictions are enforced by choose and merge
rules, and the graph is a model.  The default blocking condition is anywhere
subset blocking, maintained incrementally; pairwise, double, and Konclude's
optimised blocking are available for the shapes that need them
\cite{horrocks1999practical,steigmiller2016thesis}.  Classification follows
the model-based approach of HermiT \cite{glimm2012novel}.  The first phase tests
each class $A$ for satisfiability and takes the root label of the model as the
candidate subsumers of $A$; a node of an earlier model that already carries $A$
serves as a witness without a new test.  Told subsumers are accepted without a
test \cite{baader1994empirical}.  The second phase tests $A\sqcap\neg B$ for each
remaining candidate $B$, and every satisfiable test is a further model of $A$
that removes the candidates absent from its root label.  The tests of distinct
classes run in parallel over a shared clause index.  A separate non-branching
saturation classifier certifies its answer for nominal-free inputs with inverse
roles and defers otherwise.

The Konclude port is the largest component of KM (Section~\ref{sec:overview}). It comprises the completion
algorithm in 36 units, the approximation saturation in 12, Konclude's process
data model, caches, the classifier, and the role-chain automata
preprocessor, with Konclude's pointers replaced by arena indices.  The bridge
rebuilds Konclude's concept structures from the frontend's source axioms.
Named left sides become lazily unfolded definitions \cite{horrocks2000axioms};
structural left sides go to a triggered binary absorber that moves fully
triggerable head complements into the antecedent and otherwise unfolds a
partial trigger \cite{hudek2006binary,steigmiller2014absorption}; and
preprocessing extracts disjuncts common to all branches of a disjunction and
rewrites universal and existential restrictions over complex roles into
automaton states \cite{horrocks2006sroiq}.  The completion rules are those of
Konclude: conjunction, disjunction with alternatives ordered by Konclude's
heuristics, existential and universal successor rules, triggered implication,
at-least and at-most with merging, functionality, nominals and values, self
restrictions, automaton transitions, and datatype rules.  Backtracking restores
the complete state of an alternative from copy-on-write branch epochs, and the per-individual node table is itself stored in shared 256-entry
chunks, so a branch copies only the chunk it writes.  Before any completion,
Konclude's approximation saturation runs once over the terminology with one
node per concept \cite{steigmiller2014coupling,steigmiller2015payg}.  A subject
whose saturation node is clashed is unsatisfiable, a subject whose node is
complete and sufficient has an exact subsumer row, and only certified labels
enter the trusted subsumer sets.  Completion then reuses saturation results
for successor expansion, blocking, and cached absorption.

The remaining subjects are classified by Konclude's known- and
possible-subsumer (KPSet) procedure.  Known pairs are seeded from the told and
definitional closure of the source axioms.  In the first phase, a completion
probe for each remaining subject either opens no disjunctive choice, in which
case the named positive labels of the root are its exact subsumers, or
reports deterministic and possible subsumers, with every class absent from the
model refuted.  After a barrier over all subjects, the possible subsumptions
are pruned by propagation up and down the known hierarchy.  In the second
phase, each remaining candidate $c$ of a subject $s$ is accepted as a
deterministic subsumer, refuted by a countermodel from an earlier probe or by
transitivity through a known superclass, or decided by a probe of
$s\sqcap\neg c$.  Konclude's pseudo-model precheck is ported but not used to
refute candidates, because it produced a false refutation on one ORE ontology.
The bridge returns a complete taxonomy or defers.  It defers when an input
guard fails, when any clause cannot be encoded, when an ABox consistency
check is undecided, when a probe exhausts its budget after its retries, or
when any subject is left without a verdict.  It also defers on a condition
specific to KM, which we call poisoning: a backtrack that cannot restore the
failed alternative's state marks the affected nodes, after which a clash-free
fixpoint is not accepted as a model, so satisfiability and read-off verdicts
become unknown while unsatisfiability verdicts remain sound.

\begin{figure*}[t]
\centering
\begin{minipage}[t]{0.48\textwidth}
\centering
\begin{tikzpicture}[flow, node distance=3.2mm and 6mm]
\node[step=52mm] (a) {hypertableau clauses; clause index};
\node[test=52mm, below=of a] (b) {ontology consistent?};
\node[step=52mm, below=of b] (c) {phase 1 per class $A$: test $A$ or reuse a witness};
\node[step=52mm, below=of c] (d) {expand: hyperresolution, nominal merges, existentials on unblocked nodes};
\node[test=52mm, below=of d] (e) {open disjunction?};
\node[step=52mm, below=of e] (f) {branch; on a clash backjump or assert the negation};
\node[step=52mm, below=of f] (g) {model: root label gives candidates; accept told subsumers};
\node[step=52mm, below=of g] (h) {phase 2: test $A\sqcap\neg B$; models prune candidates};
\node[term=52mm, below=of h] (out) {taxonomy, closed transitively};
\draw[->] (a) -- (b); \draw[->] (b) -- node[lab,right]{yes} (c); \draw[->] (c) -- (d); \draw[->] (d) -- (e);
\draw[->] (e) -- node[lab,right]{yes} (f);
\draw[->] (f.east) -- ++(0.25,0) |- (d.east);
\draw[->] (e.west) -- ++(-0.25,0) |- node[lab,pos=0.25,left]{no} (g.west);
\draw[->] (g) -- (h); \draw[->] (h) -- (out);
\draw[->] (b.east) -- ++(0.55,0) |- node[lab,pos=0.05,above]{no} (out.east);
\end{tikzpicture}

\small (a) Native hypertableau
\end{minipage}\hfill
\begin{minipage}[t]{0.5\textwidth}
\centering
\begin{tikzpicture}[flow, node distance=3.2mm and 6mm]
\node[step=56mm] (a) {source axioms and hypertableau clauses};
\node[test=56mm, below=of a] (b) {guards pass and every clause encoded?};
\node[step=56mm, below=of b] (c) {lazy unfolding, binary absorption, common disjuncts, role-chain automata};
\node[step=56mm, below=of c] (d) {seed known pairs; approximation saturation answers certified subjects};
\node[step=56mm, below=of d] (e) {phase 1: probe each remaining subject; read off or report possible subsumers};
\node[step=56mm, below=of e] (f) {prune possible subsumptions};
\node[step=56mm, below=of f] (g) {phase 2: accept, refute by countermodel, or probe $s\sqcap\neg c$};
\node[test=56mm, below=of g] (h) {every subject decided?};
\node[term=56mm, below=of h] (out) {complete taxonomy and consistency};
\node[term=11mm, right=4mm of e] (x) {defer};
\draw[->] (a) -- (b); \draw[->] (b) -- node[lab,right]{yes} (c); \draw[->] (c) -- (d); \draw[->] (d) -- (e); \draw[->] (e) -- (f); \draw[->] (f) -- (g); \draw[->] (g) -- (h);
\draw[->] (h) -- node[lab,right]{yes} (out);
\draw[->] (b.east) -| node[lab,pos=0.2,above]{no} (x.north);
\draw[->] (h.east) -| node[lab,pos=0.2,below]{no} (x.south);
\end{tikzpicture}

\small (b) Konclude port through the bridge
\end{minipage}
\caption{The two hypertableau procedures.  (a) Phase 1 builds one model per
class, phase 2 tests the remaining candidate subsumers, and every satisfiable
test prunes the candidates further.  (b) Saturation answers the subjects it
can certify, the KPSet phases decide the rest, and any undecided subject makes
the whole classification defer.}
\label{fig:flow-ht}
\end{figure*}

An ontology with DL-safe rules selects the rule route.  The frontend accepts
rules whose atoms are class, object-property, and same- or different-individual
atoms, including nominals, inverses, number restrictions, data variables that
occur in the body only, and built-ins with finitely many solutions, and it
normalises further rule forms into ordinary clauses.  A consistency
stage seeds the ABox as named nominal nodes guarded by a marker concept, so
that rule variables bind only to named individuals, fires the accepted rules,
and reports inconsistency on a clash; otherwise the consequence-based engine
classifies the TBox, which is exact because DL-safe rules range over named
individuals only and cannot change a named-class subsumption.  Small rule and
data sources may instead be ground into an equivalent rule-free source before
the frontend runs.  A further complete-answer-or-defer specialist for
$\mathcal{SHOIQ}$ typed ABoxes runs the native hypertableau without blocking
and declines whenever a completed model has the configuration that would
require the nominal-introduction rule of Horrocks and Sattler
\cite{horrocks2007shoiq}.

The rule shared by all of these mechanisms is complete answer or defer.  A
specialist either returns the complete named-class taxonomy or exact
consistency decision requested, together with the evidence for that result,
or it returns a typed decline with no semantic output, after which the
supervisor invokes a fallback route designed to be complete.  The fallback
may still time out, exhaust memory, or decline, in which case KM publishes no
classification and exits with a dedicated status code.
